% Options for packages loaded elsewhere
\PassOptionsToPackage{unicode}{hyperref}
\PassOptionsToPackage{hyphens}{url}
\documentclass[
  11pt,
]{article}
\usepackage{xcolor}
\usepackage[margin=1in]{geometry}
\usepackage{amsmath,amssymb}
\setcounter{secnumdepth}{-\maxdimen} % remove section numbering
\usepackage{iftex}
\ifPDFTeX
  \usepackage[T1]{fontenc}
  \usepackage[utf8]{inputenc}
  \usepackage{textcomp} % provide euro and other symbols
\else % if luatex or xetex
  \usepackage{unicode-math} % this also loads fontspec
  \defaultfontfeatures{Scale=MatchLowercase}
  \defaultfontfeatures[\rmfamily]{Ligatures=TeX,Scale=1}
\fi
\usepackage{lmodern}
\ifPDFTeX\else
  % xetex/luatex font selection
  \setmainfont[]{Times New Roman}
\fi
% Use upquote if available, for straight quotes in verbatim environments
\IfFileExists{upquote.sty}{\usepackage{upquote}}{}
\IfFileExists{microtype.sty}{% use microtype if available
  \usepackage[]{microtype}
  \UseMicrotypeSet[protrusion]{basicmath} % disable protrusion for tt fonts
}{}
\makeatletter
\@ifundefined{KOMAClassName}{% if non-KOMA class
  \IfFileExists{parskip.sty}{%
    \usepackage{parskip}
  }{% else
    \setlength{\parindent}{0pt}
    \setlength{\parskip}{6pt plus 2pt minus 1pt}}
}{% if KOMA class
  \KOMAoptions{parskip=half}}
\makeatother
\setlength{\emergencystretch}{3em} % prevent overfull lines
\providecommand{\tightlist}{%
  \setlength{\itemsep}{0pt}\setlength{\parskip}{0pt}}
\usepackage{bookmark}
\IfFileExists{xurl.sty}{\usepackage{xurl}}{} % add URL line breaks if available
\urlstyle{same}
\hypersetup{
  hidelinks,
  pdfcreator={LaTeX via pandoc}}

\author{}
\date{}

\begin{document}

\section{Response to Reviewers}\label{response-to-reviewers}

\textbf{Manuscript:} \emph{Proxy-Based Diagnostics of Quantum Oracle
Sketching Robustness for Non-IID Sensor and Telemetry Streams}\\
\textbf{Journal:} \emph{Sensors} (MDPI) · \textbf{Manuscript ID:}
sensors-4470240

Dear Editor,

Thank you for the opportunity to submit a revised version of our
manuscript, \emph{Proxy-Based Diagnostics of Quantum Oracle Sketching
Robustness for Non-IID Sensor and Telemetry Streams} (Manuscript ID:
sensors-4470240), to \emph{Sensors}. We appreciate the time and effort
that you and the reviewers have dedicated to the paper, and we are
grateful for the constructive comments, which have improved it
considerably.

Both reviews converge on the same central point, and we agree with it.
The previous version, although hedged internally, still read as evidence
about quantum performance, whereas what the work can honestly support is
a proxy-based diagnostic and hypothesis-generation framework for
correlation-sensitive streaming analysis. The revision repositions the
manuscript accordingly, in the title, abstract, statistics, figures and
conclusions, without removing any experiment or altering any measured
result. All stochastic outputs regenerate identically from the same
seeded pipeline, which we verified before beginning the re-analysis. In
accordance with the Academic Editor's request that preprints not be
cited, the Zhao et al.~(2026) preprint has been removed from the
manuscript and the reference list, which now contains 54 references; no
other work was substituted for it, and the affected statements were
revised so that each remaining reference supports only what it actually
establishes.

The changes are marked in the highlighted version submitted with this
letter. A point-by-point response to the comments follows.

\begin{center}\rule{0.5\linewidth}{0.5pt}\end{center}

\subsubsection{Comments from Reviewer 1}\label{comments-from-reviewer-1}

\textbf{Comment 1:} \emph{I strongly recommend to reduce the emphasis on
quantum advantage and instead presenting the τ,r landscape primarily as
a diagnostic framework to identify scenarios where future quantum
implementations is most promising.}

\textbf{Response:} Thank you for this recommendation, which we have
adopted as the organising principle of the revision. The manuscript is
now presented as a diagnostic and hypothesis-generation framework rather
than as evidence of quantum advantage. The title now leads with
``Proxy-Based Diagnostics''; the Introduction states that the paper
``neither demonstrates nor claims a realised quantum-classical
separation''; and the Conclusions identify the (tau, r) landscape as the
primary contribution, with hypothesis generation as the secondary one.
Section 7.4 states three falsifiable hypotheses, H1 to H3, each with an
explicit refutation criterion, so that the question of where future
implementations are most promising becomes concrete and testable.
Throughout the paper, ``quantum-favourable'' has been replaced by
``proxy-favourable'' or ``hypothesized favourable''.

\textbf{Comment 2:} \emph{Please clarify the relationship between the
heuristic proxy model and the original theoretical framework.}

\textbf{Response:} Thank you for raising this. A dedicated subsection
now draws that boundary explicitly. Section 3.3 separates five things:
the published quantum-streaming results that provide the general
context, together with the statement that these are problem-specific and
do not establish the operational tau, r, T\_eff or accuracy model used
here; the assumptions this paper adopts; the heuristic constructions it
introduces, namely the effective-sample correction in T\_eff and the
VC-type task rate; the memory curves, retained only as illustrative
scaling references; and the scope limitations. Table 1 records the
provenance and epistemic status of every model-level quantity, with a
column stating what is not theoretically guaranteed. The relationship is
therefore presented as one of motivation and analogy rather than
derivation, and the manuscript states that it adds no quantum theorem.

\textbf{Comment 3:} \emph{The analytic Markov τ proxy differs
substantially from the empirical measurements. I recommend to discuss
this discrepancy in greater detail and explain whether it reflects
limitations of the proxy model or the synthetic data generation
process.}

\textbf{Response:} We agree that this deserved fuller treatment, and
Section 7.3 is now devoted to it. For Markov switching the two
quantities answer different questions: each bit's marginal
autocorrelation decays as rho\^{}k and crosses 1/e at about 2.8 lags at
rho = 0.7, whereas the joint n-bit chain requires on the order of n log
n / (1 - rho) steps to mix. Both values are correct answers to different
questions, so the divergence is partly expected behaviour and partly a
limitation of summarising dependence in a single scalar. It is not an
artefact of the data generation: the generators reproduce their intended
statistics, and the gap persists under all four alternative tau
estimators examined in Experiment 9, which isolates the estimator family
as the cause. The subsection also sets out the practical consequences,
namely that the empirical tau serves as the operational input, the
analytic tau as a conservative envelope, and their ratio as a diagnostic
for long-memory content.

\textbf{Comment 4:} \emph{I recommend to validate the proposed framework
on additional real-world datasets. It will help to better demonstrate
its generalizability.}

\textbf{Response:} We agree that this is the most valuable extension,
and we are explicit that it has not been carried out in the present
revision. Adding datasets would have required a new experimental
campaign, and we preferred to keep this revision re-analysis based so
that every published number remains reproducible from the pinned seeds.
What we have added, from the same raw data, is an encoder and resolution
sensitivity analysis in Section 6.8: the two streams were re-encoded
across ten combinations of sampling resolution and bucket width, and the
ordering of NYC Taxi above Machine Temperature holds in all ten. The
reported placement is therefore not an artefact of one encoder setting.
The manuscript now states that the real-stream experiment is
illustrative rather than statistically representative, the Threats
section records the two-dataset limitation, and Future Work item (iii)
commits to network telemetry, cybersecurity event streams, financial
series and IoT fleets, with hypothesis H1 defining the protocol by which
generalisability will be tested.

\textbf{Comment 5:} \emph{Please include a more comprehensive
sensitivity analysis of the selected proxy parameters, particularly C
and δ, by this you will show how they influence the reported results.}

\textbf{Response:} Agree, and we have added this quantitatively. Because
the proxy is a closed-form function of the measured tau and r, it can be
re-evaluated on the same per-seed values across C in \{1, 1.5, 2, 2.5\}
crossed with delta in \{0.01, 0.05, 0.1\}. The new Table 8 reports the
resulting movement of the proxy-baseline crossover. The ordering of the
regimes is invariant, since C and delta enter only through the monotone
factor C\^{}2 log(1/delta), but the crossover location moves from about
rho* = 0.38 at C = 2.5 and delta = 0.01 to beyond the sweep range at C =
1, with rho* about 0.82 at the operating point. The text now instructs
that the landscape be read ordinally, the abstract carries this
conditionality, and Section 5.3 states that C = 2 is a conventional
order-unity choice not fixed by any constant in the cited literature.

\textbf{Comment 6:} \emph{The comparison between deterministic quantum
proxies and stochastic classical baselines should be presented more
cautiously. The current statistical analysis may unintentionally suggest
a comparison between implemented algorithms rather than between
empirical results and theoretical proxies.}

\textbf{Response:} This is an important point and we have rebuilt the
statistical presentation around it. Table 4 now reports Hodges-Lehmann
estimates of the paired difference with exact signed-rank 95\% intervals
and matched-pairs rank-biserial correlations, and the p-values are
retained only as descriptive companions. A dedicated paragraph in
Section 6.5 explains that the pairing is between a stochastic learner
and a model prediction, that the proxy's seed-to-seed variation comes
solely from the per-seed tau estimate, and that the rank-biserial value
here indexes sign concordance across seeds rather than magnitude. The
same paragraph states that the tests ``do not, and cannot, establish
that one algorithm outperforms another, and no quantum algorithm is
implemented anywhere in this comparison''. We also disclose that none of
the five comparisons survives Holm correction at n = 10, that the
interval at rho = 0.48 is not robust to a single step of the
integer-valued tau estimator, and that the crossover location is a
deterministic solution whose uncertainty derives from the estimator and
the constants rather than from the tests.

\textbf{Comment 7:} \emph{I recommend to provide a stronger
justification for the selected classical baselines and discussing if
additional streaming methods can offer a more informative comparison.}

\textbf{Response:} Thank you for this suggestion. Section 5.2 now
justifies the choice explicitly. The three baselines span complementary
axes of bounded-memory streaming learning, namely gradient-based (online
SGD), variance-reduced gradient-based (averaged SGD) and non-gradient
hash-count-based (Count-Min), while sharing the single-pass,
fixed-memory regime that the proxy models. We also explain why
drift-adaptive methods such as ADWIN-equipped learners and adaptive
random forests are deliberately excluded from the primary comparison:
they respond to non-stationarity by discarding history, which would
confound the correlation-budget question the proxy isolates.
Benchmarking them on the same landscape is named as an extension in
Future Work.

\textbf{Comment 8:} \emph{Please explain more clearly why the proposed
operational proxies, τ and r, are sufficient to characterize the
correlation properties of the studied data streams, and discuss any
limitations of relying only on these two measures.}

\textbf{Response:} We agree, and the revision now treats two-scalar
sufficiency as an assumption rather than a fact. Section 3.3 lists it
under Assumptions as a modelling choice of this paper, and Section 7.3
shows its failure mode: for non-mixing, long-range-dependent processes
there is no single scalar on which the analytic and empirical estimators
could agree, so a single-number tau is intrinsically lossy in that
regime. Two mitigations are stated. The ratio of analytic to empirical
tau is itself a useful third diagnostic that flags long-memory content,
and multi-scale refinements of tau, using integrated-autocorrelation
estimates at several horizons, are named as follow-on work. The Threats
section records the residual risk under construct validity.

\textbf{Comment 9:} \emph{I recommend to expand the discussion of the
assumptions behind each synthetic correlation regime and to please
explain how closely these settings represent practical sensor and
telemetry environments.}

\textbf{Response:} Thank you for this suggestion. Section 4 already
carried per-regime provenance notes, and the summary in Section 4.6 now
adds the realism mapping. Markov switching idealises short-range state
persistence in sensor readouts; seasonal drift represents the diurnal
and weekly cycles of telemetry and demand series; burst repetition
represents the duplicate-heavy event bursts of monitoring and logging
pipelines; and long-range dependence represents the self-similar
behaviour first measured on Ethernet traffic. We state plainly that
these are idealisations, that real streams mix the mechanisms, as the
two NAB cases show, and that the mapping to any concrete deployment is
qualitative.

\textbf{Comment 10:} \emph{The manuscript would benefit from a dedicated
discussion of threats to validity, include the assumptions introduced by
the heuristic proxy model, synthetic data generation and feature
binarization.}

\textbf{Response:} Agree. This is now a dedicated section. Section 8,
``Threats to Validity'', is organised by construct, internal, external
and statistical-conclusion validity, and covers the items you list. It
addresses the heuristic proxy constructs, with pointers to the
provenance table; the selection of synthetic generators; feature
binarisation, noting that arcsine shrinkage makes the real-stream tau
and r estimates lower bounds on the correlation of the underlying
continuous process; the design of the deterministic-versus-stochastic
comparison; dataset selection; the definition of the target, including
the full-series threshold look-ahead raised in review; the fixed
hyperparameters; and the ten-seed budget.

\textbf{Comment 11:} \emph{I recommend to discuss the computational
complexity of estimating τ and r and the scalability of the proposed
framework for high-volume streaming systems.}

\textbf{Response:} Thank you for pointing this out. Section 5.1 now
contains a paragraph on estimator cost and scalability. The tau
estimator requires the bitwise autocorrelation up to lag K, which costs
O(nTK) time and O(nK) memory using a ring buffer per feature. The r
estimator scans a forward window of length W, costing O(TW) comparisons,
or expected O(T) with hashing, and O(Wn) memory. Both are linear in T,
independent of the alphabet size, single-pass friendly and parallel
across features. On the longest stream studied, T = 22,695, they
complete in well under a second, so the diagnostic adds negligible
overhead in high-volume deployments.

\textbf{Comment 12:} \emph{Some figures contain many overlapping curves
and annotations, I recommend to simplify the visual presentation and
enlarge labels to improve readability.}

\textbf{Response:} Thank you for this observation. All nine figures were
re-audited and several were regenerated. The legend entry ``Proxy-model
advantage'' in Figure 5 has been renamed ``Proxy-estimated separation'',
and stray typographic characters were removed from five panel titles.
In-image labels that referred to qubit counts and to quantum
order-notation were replaced with neutral proxy labels in Figures 1, 2
and 6. The synthetic markers in Figure 8(d) are now identified as
schematic orientation marks, and the caption gives the numeric stream
coordinates instead of relying on visual adjacency. Confidence-interval
statements in captions were scoped to the stochastic panels only, and
four tables that previously overran the text width were reformatted. All
regenerations are rendering changes only; the underlying numerical
outputs are unchanged.

\textbf{Comment 13:} \emph{The discussion could better distinguish
theoretical quantum memory advantages from the practical observations
obtained through the proxy-based evaluation.}

\textbf{Response:} We agree, and this distinction is now enforced
wherever the memory curves appear. Experiment 6 opens with the statement
that no memory lower bound is measured empirically in that experiment
and that no result constitutes an empirical memory floor for online SGD,
averaged SGD or Count-Min. Section 3.3, under ``Memory references'',
records that the published quantum-streaming papers do not establish the
O(n) and Omega(sqrt(N)) pair for the diagnostic task studied here, and
that the pair is retained solely as a paper-defined illustrative scaling
reference. Table 1 lists that pair as ``introduced here (illustrative
scaling reference)'', with ``a proven separation for this diagnostic
task'' in the not-guaranteed column. The caption of Figure 6 labels
panel (a) as an illustrative memory-scaling reference with no measured
content, and Section 7.1 closes its three structural factors with the
statement that all three are properties of the proxy model or its
illustrative references and that none is a measured quantum result. The
statement about n \textgreater= 20 is marked as an extrapolation beyond
the tested range.

\textbf{Comment 14:} \emph{Please discuss the applicability of the
proposed framework to other streaming domains beyond sensor telemetry,
such as network traffic, financial streams, or cybersecurity
monitoring.}

\textbf{Response:} Thank you for this suggestion. Table 9 now maps
network backbone traffic, financial returns, sensor and IoT streams, and
concept-drifted streams onto the landscape using published measurement
ranges, with the values recomputed from the paper's own formula, the
verdicts given as ranges, and the estimator convention stated. Future
Work item (iii) commits to the domains you name: network telemetry in
the form of flow and latency series, cybersecurity event streams such as
intrusion and log-anomaly feeds, financial tick and returns series, and
large-scale IoT monitoring fleets.

\textbf{Comment 15:} \emph{strengthen the conclusion, for this emphasize
the practical contributions of the study, clearly summarize its
limitations, and please outline the most important directions for future
research.}

\textbf{Response:} Agree. The Conclusions have been rewritten to this
structure. A paragraph on the practical contribution presents the
landscape as a low-cost pre-screening step: estimate tau and r, place
the stream, and read off whether correlation alone erodes the effective
budget below the level needed for rare-event learning, as it does for
the Machine Temperature stream at T\_eff = 52. A second paragraph on
limitations and required validation points to the Threats section and
names the decisive missing step, namely an implemented or simulated
sketching pipeline evaluated at coordinates the landscape predicts to be
favourable and unfavourable. The research-question answers are restated
in effect-size terms with the estimator convention made explicit in each
case.

\subsubsection{Comments from Reviewer 2}\label{comments-from-reviewer-2}

\textbf{Comment 1:} \emph{The key issue is the relationship between τ
and r and the quantities in the framework of quantum oracle-sketching.
\ldots{} They have to strengthen the theoretical basis for the chosen
proxy construction and distinguish the results from the original from
the assumptions that are used in this paper.}

\textbf{Response:} Thank you for identifying this as the key issue.
Section 3.3 now addresses it directly and states the boundary rather
than assuming an equivalence. The section has five labelled parts. The
first sets out the published quantum-streaming results that motivate the
study, namely a polynomial quantum space advantage for one-pass triangle
counting, an exponential quantum space advantage for approximating
maximum directed cut, and a black-box route from certain classical
sketches to quantum streaming algorithms, together with the statement
that these are problem-specific results that do not establish the
operational tau, r, T\_eff or accuracy model used here and that this
paper adds no quantum theorem. The remaining parts cover the assumptions
this paper adopts, the heuristic constructions it introduces, the memory
curves retained as illustrative scaling references, and the scope
limitations, which state that tau and r are measurement conventions and
are not formally equivalent to a universal mixing time, an
effective-sample count or a quantum sample-overhead parameter. Table 1
gives the provenance of each quantity with an explicit not-guaranteed
column. In direct answer to your question, tau and r are operational
summaries defined in this paper for the proxy-based diagnostic; the
manuscript neither implements the framework nor claims to prove or
reproduce any quantum result.

\textbf{Comment 2:} \emph{The terminology of quantum performance should
be moderated throughout the paper. \ldots{} we will refer to them as
proxy-based predictions instead of the actual quantum performance. This
is particularly relevant in the Abstract, Results, Discussion, figure
captions and Conclusion.}

\textbf{Response:} Agree, and a systematic terminology pass was made
over each of the surfaces you list. The term ``quantum-favourable'' no
longer appears and has been replaced by ``proxy-favourable'' or
``hypothesized favourable''. Constructions such as ``quantum performance
proxy crosses below classical performance'' have been rewritten, for
example as ``the accuracy proxy computed with the analytic tau lies
below the measured classical baseline''. The relevant column of Table 6
is now headed ``Proxy hypothesis''. Because captions alone cannot
correct text that is part of an image, Figures 1, 2 and 6 were
regenerated to remove in-image labels referring to qubit counts and
quantum order-notation. Statements that no quantum algorithm is
implemented now appear in the Abstract, in Section 3.3, in four results
subsections, in the Threats section and in the Conclusions.

\textbf{Comment 3:} \emph{The statistical comparison between classical
baselines and the proxy model is still not resolved. Our Wilcoxon tests
compare stochastic classical results for seeds to a proxy value that is
deterministic for a given value of parameters. \ldots{} We believe the
magnitude of the classical minus proxy difference and the uncertainty
may be a better gauge than simply p-values.}

\textbf{Response:} We agree, and have adopted this in full, so that
magnitude and uncertainty now lead and the tests are re-scoped to what
they can support. For each sweep point, Table 4 reports the
Hodges-Lehmann estimate of the paired difference with its exact
signed-rank 95\% interval, for example +0.042 with interval {[}+0.011,
+0.087{]} at rho = 0 and -0.046 with interval {[}-0.079, -0.003{]} at
rho = 0.88, together with the rank-biserial correlation, which we
interpret as sign concordance across seeds because the comparison is
against a near-constant reference and all magnitude information resides
in the difference estimate. The p-values are retained only as
descriptive companions. We disclose that none survives Holm correction
at n = 10 and why, and we disclose the sensitivity of the integer-valued
tau estimator, where one lag step corresponds to 0.012 in accuracy and
is sufficient to move the interval at rho = 0.48 across zero. The
interpretation paragraph states that the crossover is the deterministic
solution of the equation between the proxy accuracy and the mean
classical accuracy, with uncertainty inherited from the estimator and
the constants rather than from the tests. The full re-analysis, which
first reproduced the published per-seed results to six significant
digits, is included with the revision.

\textbf{Comment 4:} \emph{The huge discrepancy between the analytic and
empirical refreshing time is one of the most important findings in the
paper, especially for the long-range dependence regime. \ldots{} The
authors should discuss this in more detail and discuss its implications
for the interpretation of the central (τ,r) sensitivity landscape.}

\textbf{Response:} We agree that this is one of the paper's more
important findings and have promoted it accordingly. The new Section 7.3
gives the mechanism: the analytic proxies target model-level worst-case
scales, that is joint-chain mixing and tail-dominated integrated
autocorrelation, whereas the empirical estimator measures the short-lag
marginal decorrelation of the binarised stream, and for a non-mixing
process there is no scalar on which the two could agree. The
consequences for the landscape are carried through the paper. Verdicts
are now bound to their estimator convention wherever they appear, and we
state openly that the burst and long-memory extremes exchange places
between the analytic and empirical conventions, because burst saturates
through its measured duplication while the high empirical value for long
memory is a short-lag artefact. This appears in Experiment 9, in Section
7.3, in the answers to RQ2 and RQ3, and in the captions of Figures 3 and
7. The abstract now says that the two estimates diverge, rather than
that one overestimates the other, and adds that they answer different
questions.

\textbf{Comment 5:} \emph{The classical-memory comparison needs to be
taken with care. The classical lower bound curve \ldots{} is not
actually an empirical lower bound for the SGD, averaged-SGD, or
Count-Min tools used in our experiments. \ldots{} we should make it more
prominent in the dimension-scaling discussion.}

\textbf{Response:} Agree, and this is now made prominent in exactly that
discussion. The curves are no longer presented as cited bounds at all.
Experiment 6 opens with the statement that no memory lower bound is
measured empirically in the experiment, that the curves in panel (a) are
illustrative quantities, and that no result constitutes an empirical
memory floor for online SGD, averaged SGD or Count-Min. The caption of
Figure 6 labels panel (a) as an illustrative memory-scaling reference
that is paper-defined and contains no measured content, and states that
the exponential divergence of the curves is a visual scaling contrast
rather than an established separation. Section 3.3 records under
``Memory references'' that the published quantum-streaming papers do not
establish this pair of curves for the diagnostic task studied here and
that they are retained solely as a paper-defined illustrative reference.
Table 1 carries the same status, and Section 7.1 states that none of its
structural factors is a measured quantum result.

\textbf{Comment 6:} \emph{The real-stream analysis is useful as a sanity
check but it is limited to two NAB datasets. The conclusions on the
sensor and telemetry applications should therefore be conservative.}

\textbf{Response:} We agree, and the conclusions drawn from the real
streams are now explicitly conservative. Section 6.8 states that the
experiment is illustrative rather than statistically representative, the
external-validity paragraph records that generalisation is unverified,
and the abstract claims placement only. In addition, a new encoder and
resolution sensitivity table strengthens what can legitimately be said
from two streams: across ten re-encodings of the same raw series, using
three sampling resolutions and three bucket widths, the ordering of NYC
Taxi above Machine Temperature holds in every configuration, so the
reported placement is at least stable with respect to the encoder.
Broader corpora are committed to in Future Work, with a concrete test
protocol given as hypothesis H1.

\textbf{Comment 7:} \emph{the target defined based on a percentile of
the full time series should also be reconsidered in the context of
streaming analysis, since using the full series to set the target
threshold can take into account future observations. A threshold set
from a first calibration interval or updated only from previously
observed samples would provide a more rigorous prequential design.}

\textbf{Response:} Thank you for this observation, which we consider
well taken, and which we now disclose explicitly rather than leave
implicit. Section 6.8 states that the threshold is computed from the
full series and therefore uses information that would not be available
prequentially at each step. It also states that the feature stream and
the tau and r estimates involve no such look-ahead, so the landscape
placement, which is the central output for these streams, is unaffected,
and the leakage touches only the classical baseline reference numbers.
We further state that a threshold calibrated on an initial interval, or
updated only from previously observed samples, is the more rigorous
prequential design, and we adopt it as future work. The Threats section
records this as a residual risk under internal validity. We chose
disclosure over re-running the baselines in this revision in order to
keep every published number reproducible, and we would be glad to re-run
them with a calibration-interval threshold in a further round if you
prefer.

\textbf{Comment 8:} \emph{The interpretation of the NYC Taxi Count-Min
result also needs to be qualified. The six-bit representation gives a
relatively small number of possible feature vectors and hence memorising
the same feature patterns is relatively easy. \ldots{} I would like to
give more weight to the implications when we think about the apparent
agreement between the Count-Min performance and our proxy prediction.}

\textbf{Response:} Agree, and the qualification now carries weight in
three places. The dedicated paragraph explaining why Count-Min performs
well is retained, and the interpretation now states plainly that the
favourable-side separation between the two streams rests on Count-Min
alone and that the hashed linear learners do not corroborate it. We
additionally report that both SGD variants fall below the majority
baseline on accuracy and that the balanced accuracy of averaged SGD
straddles chance, so the apparent agreement between Count-Min and the
proxy prediction is explicitly not read as validation. Section 7.4 lists
representation-specific effects of this kind among the phenomena the
landscape cannot anticipate, and, going further in the direction you
indicate, the same section concedes that, read strictly as an accuracy
projection, the proxy is falsified on the Machine Temperature stream,
where 0.500 is predicted and 0.72 to 0.80 observed, with the rare-class
diagnostic reading offered as interpretation rather than confirmation.
Table 6 now carries the ordering-to-ordering intervals, with Count-Min
F1 at 0.644 plus or minus 0.062 and a per-ordering range of 0.43 to
0.72, which tempers the headline number.

\textbf{Comment 9:} \emph{In particular, the target function ablation
should be highlighted because it shows that our classical performance
can change dramatically with the classification task in question and
that our proxy is in fact target independent. A correlation-based proxy
that is insensitive to the target function cannot necessarily predict
the performance of algorithms for different kinds of decision problems.}

\textbf{Response:} We agree, and this is now highlighted and stated as a
scope boundary in your own terms. Section 6.9 states that classical
learner performance is strongly target-dependent while the proxy depends
only on tau, r and T, and is target-independent by construction, so that
the proxy cannot universally predict downstream classifier performance
and instead characterises the correlation structure of the stream, that
is the effective sample budget available to any bounded-memory learner,
rather than task-specific learnability. Section 7.4 elevates this to the
list of things the landscape cannot do, adding the empirical observation
that measured classical accuracy is flat while T\_eff falls fivefold, so
the budget model does not predict even the classical arm's accuracy. The
caption of Figure 9 states that the proxy does not track task-specific
learnability, and the answer to RQ2 now describes tau and r as budget
summaries rather than accuracy predictors. The protocol of the ablation,
which is prequential with a learning rate of 0.01 and T = 4096, is also
stated so that its absolute levels are not misread against the main
experiments.

\textbf{Comment 10:} \emph{the strongest contribution of the paper is
the (τ,r) sensitivity landscape proposed as a hypothesis-generating tool
\ldots{} The paper would be better off having this methodological
contribution as the main message rather than the separation of
quantum-classical performance.}

\textbf{Response:} Thank you for this closing recommendation, which we
have taken as the identity of the revised paper. The title leads with
``Proxy-Based Diagnostics''; the abstract announces a proxy-based
diagnostic and hypothesis-generation framework; the contribution list
names the landscape as the paper's central diagnostic and
hypothesis-generation artifact; Section 7.4 states the falsifiable
hypotheses it generates; and the Conclusions open by identifying the
landscape as the primary contribution and hypothesis generation as the
secondary one, followed by the sentence that the paper demonstrates no
quantum advantage and implements no quantum algorithm.

\begin{center}\rule{0.5\linewidth}{0.5pt}\end{center}

We hope that the revised manuscript now meets the standards of
\emph{Sensors}, and we look forward to hearing from you in due course.
We would be glad to respond to any further questions or comments.

Sincerely,

\textbf{AbdelMoniem Helmy}, on behalf of all authors Corresponding
author
\href{mailto:abdelmoniem.hafez@cu.edu.eg}{\nolinkurl{abdelmoniem.hafez@cu.edu.eg}}
· ORCID 0000-0001-5996-6019

\end{document}
